% Topic 1.2.07 — Stochastic Processes: Finite-Dimensional Distributions.

\section{Stochastic Processes: Finite-Dimensional Distributions}
\label{sec:1.2.07-stochastic-processes}

\begin{ticket}
Definition of a stochastic process. System of finite-dimensional distributions of a process and its properties. Sample space of a process.
\end{ticket}

\subsection*{Theory}

A stochastic process is the probabilistic analogue of a function:
instead of a single random outcome, we observe a whole indexed
family of outcomes. The index set might be discrete (time-step,
labelling), continuous (real time), or multi-dimensional
(spatial); the values may be real, vector-, or set-valued. The
fundamental objects are the joint laws of \emph{finitely many}
coordinates.

\begin{definition}[Stochastic process]
\label{def:stochastic-process}
A \emph{stochastic process} (or \emph{random process}) on a
probability space $(\Omega, \mathcal{F}, \Prob)$ with index
set~$T$ is a family $X = (X_t)_{t \in T}$ of random variables
$X_t \colon \Omega \to \R$. For a fixed outcome
$\omega \in \Omega$, the function $t \mapsto X_t(\omega)$ is the
\emph{sample path} (or \emph{trajectory}) of~$X$.
\end{definition}

When $T = \N$ or $\Z$ we speak of a \emph{discrete-time} process;
when $T = [0, \infty)$ or~$\R$, a \emph{continuous-time} process.
$\R$-valued processes are the default; vector- and set-valued ones
are obtained by allowing $X_t$ to take values in $\R^d$ or some
other measurable target.

\medskip
\noindent\textbf{Path space.} Each outcome~$\omega$ produces a
function $t \mapsto X_t(\omega)$ in the function space $\R^T$.
Identifying~$\omega$ with its trajectory, we may take the
\emph{path space} $\R^T$ itself as the sample space, equipped with
the product (cylinder) $\sigma$-algebra
$\mathcal{B}(\R^T) = \bigotimes_{t \in T} \mathcal{B}(\R)$, the
smallest $\sigma$-algebra making every coordinate projection
$\pi_t \colon (x_s)_{s \in T} \mapsto x_t$ measurable. With this
convention, the law of the whole process is a single probability
measure on $(\R^T, \mathcal{B}(\R^T))$.

\medskip
\noindent\textbf{Finite-dimensional distributions.} The product
$\sigma$-algebra is generated by \emph{cylinder sets} of the form
$\{(x_s) : (x_{t_1}, \dots, x_{t_n}) \in B\}$ for finite
$\{t_1, \dots, t_n\} \subseteq T$ and Borel
$B \in \mathcal{B}(\R^n)$. The probability measure is therefore
determined by the joint distributions of finite subfamilies
$(X_{t_1}, \dots, X_{t_n})$.

\begin{definition}[Finite-dimensional distributions]
\label{def:fdd}
The \emph{system of finite-dimensional distributions (FDDs)} of a
process~$X$ is the family of probability measures
\[
  \mu_{t_1, \dots, t_n}(B) \;=\; \Prob\bigl((X_{t_1}, \dots, X_{t_n}) \in B\bigr), \qquad B \in \mathcal{B}(\R^n),
\]
indexed by finite ordered tuples $(t_1, \dots, t_n) \in T^n$ for
all $n \geq 1$.
\end{definition}

\begin{proposition}[Consistency of FDDs]
\label{prop:fdd-consistency}
The FDDs of any stochastic process satisfy two
\emph{Kolmogorov consistency conditions}:
\begin{enumerate}[label=(C\arabic*)]
  \item \emph{permutation:} for any permutation~$\pi$ of
        $\{1, \dots, n\}$ and Borel
        $B_1, \dots, B_n \in \mathcal{B}(\R)$,
        \[
          \mu_{t_{\pi(1)}, \dots, t_{\pi(n)}}(B_{\pi(1)} \times \cdots \times B_{\pi(n)}) = \mu_{t_1, \dots, t_n}(B_1 \times \cdots \times B_n);
        \]
  \item \emph{marginalisation:} for any
        $B \in \mathcal{B}(\R^n)$,
        \[
          \mu_{t_1, \dots, t_n, t_{n+1}}(B \times \R) = \mu_{t_1, \dots, t_n}(B).
        \]
\end{enumerate}
\end{proposition}
\proofof{fdd-consistency}

The converse — that any consistent family of probability measures
arises as the FDDs of some process — is the content of Kolmogorov's
existence theorem.

\begin{theorem}[Kolmogorov's existence / extension theorem]
\label{thm:kolmogorov-extension}
Let $T$ be an arbitrary index set, and suppose for every finite
tuple $(t_1, \dots, t_n) \in T^n$ a probability measure
$\mu_{t_1, \dots, t_n}$ on $\mathcal{B}(\R^n)$ is given. If the
family $\{\mu_{t_1, \dots, t_n}\}$ satisfies the consistency
conditions (C1)–(C2) of \cref{prop:fdd-consistency}, then there
exists a probability space $(\Omega, \mathcal{F}, \Prob)$ and a
stochastic process $(X_t)_{t \in T}$ on it whose FDDs coincide
with the given family.
\end{theorem}
\proofof{kolmogorov-extension}

\begin{remark}
The extension theorem produces a process on the path space
$(\R^T, \mathcal{B}(\R^T))$. This carrier space is large enough to
encode all FDDs but \emph{small} in another sense: events that
depend on uncountably many coordinates — for instance,
$\{\sup_{t \in [0, 1]} X_t \leq c\}$ for $T = [0, 1]$ — need not
lie in $\mathcal{B}(\R^T)$. Resolving this typically requires the
choice of a \emph{regular version} of the process, e.g., one with
continuous sample paths (\cite[Ch.~II,~\S~2]{Sevastyanov1982}).
\end{remark}

\medskip
\noindent\textbf{Equivalent processes.} Two processes $X$ and $Y$
on possibly different probability spaces but the same index
set~$T$ are called \emph{equivalent} (or said to have the
\emph{same law}) if their FDD systems coincide. The extension
theorem shows that the FDDs are a complete invariant on the
$\mathcal{B}(\R^T)$ level, even though equivalent processes may
differ as functions of $\omega$ (the so-called modification
issue).

\subsection*{Proofs}

\begin{proof}[\Cref{prop:fdd-consistency}]
\proofanchor{fdd-consistency}
Both conditions follow from the standard properties of joint
distributions of random variables on a fixed probability space.

\emph{(C1)} For any permutation~$\pi$, the two events
\[
  \{(X_{t_1}, \dots, X_{t_n}) \in B_1 \times \cdots \times B_n\}
  \quad \text{and} \quad
  \{(X_{t_{\pi(1)}}, \dots, X_{t_{\pi(n)}}) \in B_{\pi(1)} \times \cdots \times B_{\pi(n)}\}
\]
are equal as subsets of~$\Omega$ — both reduce to
$\{X_{t_i} \in B_i \text{ for every } i\}$ — so they have equal
probability.

\emph{(C2)} The set
$\{(X_{t_1}, \dots, X_{t_n}, X_{t_{n+1}}) \in B \times \R\}$ equals
$\{(X_{t_1}, \dots, X_{t_n}) \in B\}$, since the constraint
``$X_{t_{n+1}} \in \R$'' is automatic. Hence the corresponding
probabilities coincide.
\end{proof}

\begin{proof}[\Cref{thm:kolmogorov-extension}, sketch]
\proofanchor{kolmogorov-extension}
Let $\mathcal{C}$ be the algebra of \emph{cylinder sets} in $\R^T$:
sets of the form
$C = \{(x_s) : (x_{t_1}, \dots, x_{t_n}) \in B\}$ for some finite
$F = \{t_1, \dots, t_n\} \subseteq T$ and $B \in \mathcal{B}(\R^n)$.
Define $\nu(C) = \mu_F(B)$. Consistency (C1)–(C2) makes this
well-defined: distinct presentations of the same cylinder
correspond to permutations or to embeddings into a larger finite
index set, both of which preserve the value of $\nu$.

The map $\nu \colon \mathcal{C} \to [0, 1]$ is finitely additive
(again by reduction to a common finite index set) and
$\sigma$-additive on~$\mathcal{C}$ — the latter is the non-trivial
step, established by a compactness argument that reduces to
Tikhonov's theorem in $\R^F$ (essentially: a decreasing sequence
of cylinders with non-zero measures has a non-empty intersection).
Carathéodory's extension theorem then produces a unique
$\sigma$-additive extension $\Prob$ to $\mathcal{B}(\R^T)$, and
$X_t(\omega) = \omega_t$ has the prescribed FDDs by construction.
Full details: \cite[Ch.~II, \S~3]{Shiryaev1989}.
\end{proof}

\subsection*{Examples}

\begin{example}[Bernoulli process]
\label{ex:bernoulli-process}
Let $T = \N$ and let $(X_n)$ be i.i.d.\ $\mathrm{Bernoulli}(p)$.
The FDDs are products: for any finite
$\{n_1, \dots, n_k\} \subseteq \N$ and
$\varepsilon_i \in \{0, 1\}$,
\[
  \Prob(X_{n_1} = \varepsilon_1, \dots, X_{n_k} = \varepsilon_k) = \prod_{i=1}^{k} p^{\varepsilon_i} (1 - p)^{1 - \varepsilon_i}.
\]
Both consistency conditions are immediate (the right-hand side is
permutation-invariant up to reordering and marginalisation reduces
the product). The Kolmogorov extension theorem
(\cref{thm:kolmogorov-extension}) provides the underlying
probability space, but for this example $\Omega = \{0,1\}^{\N}$
with the product $\sigma$-algebra and product measure works
directly.
\end{example}

\begin{example}[Simple symmetric random walk]
\label{ex:simple-rw}
With $(X_n)$ as above, $p = 1/2$, set $S_0 = 0$ and
$S_n = (2 X_1 - 1) + \cdots + (2 X_n - 1)$ — the cumulative sum of
$\pm 1$ steps. The FDDs of $(S_n)_{n \geq 0}$ are determined by
those of $(X_n)$ via a finite linear map, and consistency is
inherited. The conditional law $S_{n+1} \mid S_n$ is
$\frac{1}{2}\delta_{S_n + 1} + \frac{1}{2}\delta_{S_n - 1}$ — the
simplest example of a discrete-time Markov chain
(\cref{sec:1.2.09-markov-chains}).
\end{example}

\begin{example}[Gaussian process specified by FDDs]
\label{ex:gaussian-process}
Let $T \subseteq \R$, fix a function $m \colon T \to \R$ and a
positive-semidefinite kernel $K \colon T \times T \to \R$ (i.e.,
for any finite $\{t_1, \dots, t_n\}$, the matrix
$(K(t_i, t_j))_{i, j}$ is symmetric and positive-semidefinite).
Declare $\mu_{t_1, \dots, t_n}$ to be the multivariate normal law
on $\R^n$ with mean vector $(m(t_1), \dots, m(t_n))$ and covariance
matrix $K(t_i, t_j)$. The consistency conditions follow from the
fact that marginal distributions of multivariate Gaussians are
Gaussian (with the inherited mean and covariance). Hence by
Kolmogorov's theorem there exists a stochastic process — a
\emph{Gaussian process} — with the prescribed mean and covariance
structure. The most famous instance is the \emph{Wiener process}
(Brownian motion) with $m \equiv 0$ and $K(s, t) = \min(s, t)$ on
$T = [0, \infty)$.
\end{example}

\begin{problem}
\label{prob:fdd-marginal}
Suppose a stochastic process $(X_t)_{t \in [0, 1]}$ has FDDs
$\mu_{t_1, \dots, t_n}$ that are jointly Gaussian with
$\E X_t = 0$ and $\Cov(X_s, X_t) = e^{-|s - t|}$. Verify the
consistency conditions and compute the law of $X_t$ for a single~$t$
and of $(X_s, X_t)$.
\end{problem}

\begin{proof}[Solution]
$\Cov(X_t, X_t) = e^{0} = 1$, so $X_t \sim \mathcal{N}(0, 1)$ for
every~$t$. For any pair $s \leq t$,
$(X_s, X_t)$ is jointly Gaussian with mean~$\vect{0}$ and
covariance $\begin{pmatrix} 1 & e^{-(t - s)} \\ e^{-(t - s)} & 1 \end{pmatrix}$ —
in particular, $\rho(X_s, X_t) = e^{-(t - s)}$, decaying
exponentially in the time gap. Permutation consistency: the
covariance kernel is symmetric, so permuting indices reorders both
vector entries and matrix entries consistently. Marginalisation:
the marginal of an $(n+1)$-dimensional Gaussian on the first
$n$~coordinates is the $n$-dimensional Gaussian with the
corresponding mean and covariance — this is the standard fact about
Gaussian marginals. Hence Kolmogorov's existence theorem
(\cref{thm:kolmogorov-extension}) provides the process; it is the
\emph{Ornstein–Uhlenbeck} process, the stationary Gaussian process
that surfaces in mean-square ergodicity questions
(\cref{sec:1.2.08-ergodicity}).
\end{proof}
